\section{Discussion}

This discussion is intentionally result-agnostic. It does not claim that embodied biofeedback changed behavior, nor that participants experienced the button as self-like, nor that novelty outweighed physiology. Instead, it states the interpretive commitments that should govern a later data-bearing version of the paper. The central claim is conceptual: pressability is not reducible to visibility or compliance. It is the relation produced when cultural symbolism, design invitation, bodily reach, social legibility, and embodied coupling converge on a single actionable object \citep{Plotnick2018,Gaver1991,McCarthyWright2004}.

\subsection{If physiological coupling increases pressability}

If live cardiac and respiratory coupling increase pressing, presence, or felt closeness to the button relative to sham feedback, the most defensible interpretation is not that the button ``became the self'' in any simple sense. The stronger claim would be narrower: physiological synchrony can alter how a reachable object is integrated into the participant's multisensory and action-oriented field. That reading would be consistent with prior work on embodiment, peripersonal space, and biofeedback-mediated VR experience \citep{KilteniEtAl2012,Serino2019,RockstrohEtAl2019,ChittaroEtAl2024}. Such a result would make the button interesting not because it is magical, but because it becomes a tractable site where symbolic invitation and bodily coupling can be manipulated together.

\subsection{If symbolic saturation and ergonomics dominate}

Mixed or null condition effects would still matter. A large red button in a sparse room may already be so culturally legible and perceptually loud that additional physiology changes little. Plotnick's histories, salience research, and placement studies all point in this direction: the object arrives with a substantial amount of explanatory freight before the first heartbeat-triggered pulse occurs \citep{PlotnickInterface2012,PlotnickWarfare2012,ElliotMaier2014,Li2008,Lin2021}. Under that reading, the study would show that some interface objects are already overdetermined. Embodiment may modulate pressability at the margins, while symbol and reach set the baseline.

\subsection{If novelty overwhelms condition effects}

The repository's most conceptually interesting possibility may be the simplest one: the first encounter dominates. Curiosity and novelty research suggest that first contact can reorganize attention and memory in ways that are difficult to repeat \citep{KiddHayden2015,RanganathRainer2003}. The BRB protocol is unusually explicit about this problem. Participants are invited to treat the second round as if it were the first, while the study simultaneously acknowledges that this is impossible. The paired language of oneness and secondness is therefore methodologically awkward but theoretically productive. It reframes order effects as part of the phenomenon rather than as mere contamination. If first encounter turns out to be the strongest predictor of behavior, that would not be a failed manipulation. It would be a finding about the temporal irreversibility of pressability.

\subsection{Contribution, display, and the micro-politics of the counter}

The visible press counter deserves special attention because it changes what a press is for. It does not merely confirm successful input; it frames each action as a contribution to a shared empirical record. In experience-centered terms, the participant is invited to become a co-producer of the dataset rather than a hidden source of responses \citep{McCarthyWright2004}. This may increase agency, performance consciousness, or even playful compliance. It also means that pressing behavior cannot be interpreted as purely private desire. The interaction is staged as methodologically consequential, and that staging is part of the phenomenon under study.

\subsection{Limitations and inclusive bounds}

Several limits should remain explicit even after the data arrive. First, the current repository does not yet contain the data needed for a completed empirical paper. Second, secondness is presently a bespoke measure without archived validation. Third, red cannot be treated as a universally stable sign for danger, urgency, or invitation; color-emotion mappings vary culturally, and any discussion of ``the'' meaning of a red button must remain bounded \citep{GaoEtAl2007,ElliotMaier2014}. Fourth, bodily variation matters. Arm length, handedness, mobility, prior VR familiarity, and comfort with chest-strap sensing can all shape what counts as a reachable or acceptable interaction \citep{Bonin2022,CummingsBailenson2016}. Finally, the button is an unusually saturated object. That is analytically useful, but it also means the paper should resist generalizing too quickly from one iconic artifact to all interface actions.

Taken together, these limits do not weaken the project; they define the terms under which it can remain credible. The big red button is already halfway to theory before the participant enters the headset. The task of the paper is to show which parts of that theory survive contact with actual data.
